%% kapa: exam=2 round=27 year=2016 period=2 subject=appraisal_theory subject_ko=감정평가이론 seq=1 number=1 points=40
%% source: Q-Net 감정평가사 기출문제 (cst00309) artlSeq=5185401 file=extracted/1008/5185401/제27회 감정평가사 2차시험 기출문제/제27회 감정평가사 2차 시험문제지/(2교시 시험문제지)감정평가이론.pdf
%% conversion: kordoc 4.24.51 -> tools/gen2.py (auto); review=proofread-gpt-sol; flags=none
\begin{longquestion}{【문제 1】}
지식정보사회로의 이행 등에 따라 기업가치 중 무형자산의 비중(Portion)이 상대적으로 증가하고 있다. 『감정평가 실무기준』에 규정하고 있는 계속기업가치(going concern value)의 감정평가와 관련하여 다음 물음에 답하시오. (40점) \par
1) 기업가치의 구성요소를 설명하고, 기업가치의 감정평가 시 유의사항을 설명하시오. (10점) \par
2) 기업가치의 감정평가에 관한 이론적 배경과 감정평가방법을 설명하고, 각 감정평가방법 적용 시 유의사항 및 장단점을 설명하시오. (20점) \par
3) 기업가치의 감정평가에 있어서 시산가액 조정에 대하여 설명하고, 조정된 기업가치에 대한 구성요소별 배분방법에 관해 설명하시오. (10점) \par
\end{longquestion}
